{\Large
	\noindent
	{\bf Abstract}} \\
\vspace*{0.3 cm}

\noindent
The increasing dependence of modern electricity systems on weather-sensitive renewable generation has made anomaly detection in power markets and grid operations both more important and more challenging. Traditional statistical methods often struggle to distinguish between truly abnormal behaviour and context-driven variation caused by meteorological conditions, seasonal effects, and market volatility. This thesis addresses that challenge by proposing a lightweight neuro-symbolic framework for time-series anomaly detection in the Austrian electricity domain. The framework combines a compact discrete Cellular Neural Network and Extreme Learning Machine forecasting backbone (dCeNN--ELM) with Answer Set Programming (ASP) for symbolic refinement of anomaly candidates.

The proposed system is built on an engineered Austrian hourly dataset covering the period 2017--2022 and integrates electricity-market variables, grid-operational signals, weather data, installed-capacity information, and calendar context. The neural component models expected multivariate system behaviour and generates residual-based anomaly candidates, while the symbolic component applies persistence, plausibility, and domain-consistency rules to reject weak or implausible detections. In this way, the framework moves beyond purely statistical anomaly scoring toward more interpretable and operationally meaningful anomaly analysis.

The empirical evaluation considers forecasting performance, anomaly refinement, finance-aware utility, and event-based case studies. The results show that the proposed framework does not universally outperform all baselines in pure forecasting accuracy; in particular, Ridge Regression remains a strong benchmark in the committed repository results. However, the dCeNN--ELM model achieves substantially lower training and inference cost than the LSTM baseline and is better aligned with the edge-oriented and modular design goals of the thesis. More importantly, the ASP refinement stage improves interpretability by providing explicit veto logic for rejected anomalies and supports false-alarm reduction through domain-aware reasoning. Event-level analysis further demonstrates the framework’s ability to capture prolonged price disturbances, severe load deviations, renewable-side stress episodes, and multi-signal anomaly events.

Overall, this thesis shows that a lightweight neuro-symbolic approach can provide a credible and practically useful alternative to purely statistical anomaly detection in weather-dependent electricity systems. The main contribution lies not in claiming universal predictive superiority, but in demonstrating how learning, reasoning, and downstream utility analysis can be integrated into one coherent anomaly-analysis pipeline.